\section{Physical pressure domination and complex Gaussian estimates}
\label{sec:v59-contact-details}
We give the scalar details used in the pressure-contact argument before
turning to a correlated realization with a continuous roof. This section
uses no collision norm. In particular, the pressure comparison is an
inequality between original scalar integrals, not an order relation on
an anisotropic space.

\begin{lemma}[Taking roots of a visible physical branch]
\label{lem:v59-physical-root}
Let $X$ be a centered bounded observable of a probability-preserving
system. Suppose, at a fixed real tilt $t$,
\[
 1\le\int e^{-t\cdot S_mX}\,d\nu\le C_0e^{m\mathscr P(t)}.
\]
Suppose bounded physical endpoint tests $a,d$ give, at $z=a_0+it$,
\[
 I_m(z)=A(z)e^{m\psi(z)}+E_m(z),\qquad
 |A(z)|\ge a_*>0,\qquad |E_m(z)|\le C_1\rho^m,
 \quad\rho<1,
\]
where $I_m(z)=\int a e^{iz\cdot S_mX}(d\circ T^m)d\nu$ and the
logarithm is for this \emph{centered} observable. Then
\begin{equation}\label{eq:v59-physical-root}
 \mathscr P(t)\ge0,\qquad \operatorname{Re}\psi(a_0+it)\le\mathscr P(t).
\end{equation}
The conclusion is uniform on any finite cover with common bounds and a
common tilt ball. The tests may change between members of that cover.
\end{lemma}
\begin{proof}
The first inequality in the hypothesis follows from centering and
Jensen whenever only the upper bound is given. Taking logarithms,
dividing by $m$, and passing to the limit gives $\mathscr P(t)\ge0$.
Physical domination and the reverse triangle inequality give
\[
 a_*e^{m\operatorname{Re}\psi(z)}
 \le C_0\|a\|_\infty\|d\|_\infty e^{m\mathscr P(t)}+C_1\rho^m.
\]
For nonnegative $u_m,v_m$, $(u_m+v_m)^{1/m}$ has limsup at most
the larger of their root limsups, since $u_m+v_m\le2\max(u_m,v_m)$.
Thus $e^{\operatorname{Re}\psi(z)}\le\max(e^{\mathscr P(t)},\rho)
=e^{\mathscr P(t)}$. The lower amplitude bound is used before taking
roots; the observation amplitude of another pairing need not be nonzero.
\end{proof}

\begin{lemma}[Finite witnesses near a compact zero set]
\label{lem:v59-witness-cover}
Let finitely many compact branch supports carry continuous nonnegative
dampings $\kappa_j$. Suppose every point of their zero sets has an open
neighborhood on which a physical expansion as above holds. Then some
$\delta_0>0$ has the following property: every retained branch point with
$\kappa_j\le\delta_0$ is in one of finitely many such neighborhoods.
All constants and a smaller positive tilt radius can be chosen in common.
\end{lemma}
\begin{proof}
Choose a finite subcover of the zero set. Its complement in the compact
disjoint union of branch supports is compact and contains no zero, so
$\kappa$ has a positive minimum there, unless the complement is empty.
Choose $\delta_0$ below that minimum. Take the maximum of finitely many
upper bounds, the minimum of the positive amplitude lower bounds, and
the minimum of smaller tilt radii. No continuous global choice of
witnesses is required.
\end{proof}

\begin{lemma}[Accretive Gaussian differentiation]
\label{lem:v59-accretive-gaussian}
Let $\mathcal H$ be a compact family of complex symmetric matrices
satisfying $\operatorname{Re}H\ge\lambda I$ for some $\lambda>0$.
The Gaussian
\[
 G_H(Z)=(2\pi)^{-d/2}(\det H)^{-1/2}
                         e^{-Z^{\mathsf T}H^{-1}Z/2}
\]
uses the determinant square root continued from the real positive
matrices in the convex accretive domain. There are $C,b>0$ such that,
for real $Z$ and complex symmetric $E$,
\begin{equation}\label{eq:v59-gaussian-derivatives}
 |G_H(Z)|+|\nabla_ZG_H(Z)|+
        \sup_{\|E\|\le1}|D_HG_H(Z)[E]|\le Ce^{-b|Z|^2}.
\end{equation}
The same assertion holds on all line segments from $\mathcal H$ to any
compact family of real positive matrices with a common lower bound.
\end{lemma}
\begin{proof}
Symmetry identifies the adjoint with the entrywise conjugate. The
Hermitian identity is
\begin{equation}\label{eq:v59-inverse-real-part}
 \operatorname{Re}(H^{-1})
   =(H^*)^{-1}(\operatorname{Re}H)H^{-1}
   =\overline H^{-1}(\operatorname{Re}H)H^{-1}.
\end{equation}
Indeed multiplying $(H+H^*)/2$ on these two sides gives
$((H^*)^{-1}+H^{-1})/2$. If $\|H\|\le M$, its quadratic form is
at least $\lambda M^{-2}|Z|^2$. In particular every matrix in the
accretive domain is invertible. The domain is convex, so the analytic
logarithm of its nonvanishing determinant is fixed by its real-positive
value; this defines the indicated square root unambiguously.

Differentiating gives
\[
 \nabla_ZG_H=-H^{-1}ZG_H,\qquad
 D_HG_H[E]=\tfrac12G_H
 \{Z^{\mathsf T}H^{-1}EH^{-1}Z-\operatorname{tr}(H^{-1}E)\}.
\]
The inverse and determinant factors are uniformly bounded on the compact
family. Absorb $1+|Z|+|Z|^2$ by reducing the Gaussian exponent. Line
segments retain the positive real-part bound and form a compact set,
so the same argument applies there.
\end{proof}

\paragraph{The explicit diagonal and the two covariance moduli.}
In Section~\ref{sec:v58-pressure-contact}, $h_j$ is the norm of the real
quadratic pressure gap, whereas $\omega(\delta)$ is the supremum of
the full complex Hessian mismatch. The latter controls the matrix
segment in Lemma~\ref{lem:v59-accretive-gaussian}. Taking
\[
 \delta_m=\min\{\delta_0/2,m^{-1/2}\}
\]
in \eqref{eq:v58-shape-modulus} makes its right side tend to zero:
$\omega(\delta_m)\to0$, $m^{-1/6}\to0$, and eventually
$e^{-m\delta_m/4}=e^{-\sqrt m/4}$. This is an explicit diagonal in a
proved scalar inequality; it is not a substitution of a growing band
in a fixed-band spectral theorem. No rate is assigned to $\omega$.

\paragraph{Source normalization and exact labels.}
The first $1/c$ in $\Phi$ is the restricted collision-source
normalization. The second, in $\alpha=\Phi/c$, is the product of the
count prefactor and the covariance determinant change:
$m^{-2}\sim c^2n^{-2}$ and
$g_{\Omega_R}(\sqrt cL_RV)=c^{-3}g_{D_R}(V)$.
A torus-frequency lift changes the phase by $2\pi$ times an exact
integer label, hence leaves it unchanged. None of these operations
erases a zero arithmetic class. The reference is signed until its
positive part is explicitly normalized. In
Proposition~\ref{prop:v58-reference-positivity}, the true density is
extended by zero beyond its roof support and the local integral is
uniform over every translated unit interval. That argument controls
negative reference heights, not positive physical-source heights.
